% id: 183-437
% book: 개념원리 공통수학2 · 22 명제의 역과 대우
% section: 필수·발전 예제
% figure: none
% answer: ③
% answer_source: 답지
% variant_level: 0 (원본 전사)
% 이 파일은 build.mjs 가 items.json 에서 만든다. 고칠 때는 items.json 을 고친다.
\smash{\raisebox{1.5mm}{\dmkichul{개념원리 공통수학2 183쪽 437번 · 확인체크}}}
네 조건 $p$, $q$, $r$, $s$에 대하여 세 명제 $p\longrightarrow q$, $\sim q\longrightarrow\ \sim r$, $s\longrightarrow\ \sim q$가 모두 참일 때, 다음 명제 중 항상 참이라고 할 수 \underline{없는} 것은?
\begin{choices32}{$\sim q\longrightarrow\ \sim p$}{$p\longrightarrow\ \sim s$}{$p\longrightarrow r$}{$r\longrightarrow\ \sim s$}{$s\longrightarrow\ \sim p$}\end{choices32}
